\documentclass[a4paper,11pt]{article}

\usepackage[margin=2.5cm]{geometry}
\usepackage{hyperref}
\usepackage{listings}
\usepackage{xcolor}
\usepackage{booktabs}
\usepackage{amsmath}
\usepackage{amssymb}

\hypersetup{
    colorlinks=true,
    linkcolor=blue,
    urlcolor=blue,
    citecolor=blue
}

\definecolor{codebg}{rgb}{0.95,0.95,0.95}

\lstset{
    backgroundcolor=\color{codebg},
    basicstyle=\ttfamily\small,
    breaklines=true,
    frame=single,
    rulecolor=\color{gray},
    xleftmargin=1em,
    xrightmargin=1em,
    showstringspaces=false,
    columns=flexible,
    keepspaces=true
}

\title{espol-cuboids \\ Polarizability Sweep Suite --- User Guide}
\author{}
\date{\today}

\begin{document}

\maketitle
\tableofcontents
\newpage

\section{Introduction}

This suite computes the electrostatic polarizability tensor of dielectric
cuboids using three independent numerical methods:
%
\begin{itemize}
    \item \textbf{MoM 3D} — 3D surface integral equation (Method of Moments)
          for rectangular cuboids of arbitrary aspect ratio.
    \item \textbf{MoM 2D} — 2D boundary integral equation for infinitely long
          prisms with rectangular cross-section.
    \item \textbf{DDA} — Discrete Dipole Approximation (volume integral
          equation) on a cubic lattice inside the cuboid.
\end{itemize}
%
All three methods output the normalised polarizability tensor and effective
depolarization factors \(N_x, N_y, N_z\) for each shape and permittivity.
The MoM and DDA solvers also perform convergence extrapolation to the
continuum limit.

\section{Installation}

\begin{lstlisting}[language=,caption={Clone and install}]
git clone https://github.com/atipaqo/espol-cuboids.git
cd espol-cuboids
pip install -e .
\end{lstlisting}

Requirements: Python $\ge$ 3.9, NumPy $\ge$ 1.20, SciPy $\ge$ 1.7.

All scripts and library code are accessible after installation.  The
\texttt{espol\_cuboids} package can be imported anywhere:

\begin{lstlisting}[language=Python,caption={Library usage example}]
from espol_cuboids import (
    gen_mesh_3d_rect, comp_polarizability_3d, Mesh
)
r, n, ds, ne, V = gen_mesh_3d_rect(1.0, 1.0, 1.0, nd=30, octant=True)
P = comp_polarizability_3d(Mesh(r, n, ds, V=V, er=10.0, octant=True))
print(P.xx)  # ~2.49
\end{lstlisting}

\section{Project Structure}

\begin{lstlisting}[language=,caption={Key files and directories}]
espol-cuboids/
  pyproject.toml

  src/espol_cuboids/              # Importable package
    __init__.py                   # Public API
    vector_analysis.py            # V3, D3, Ori, Mesh data structures
    dda_core.py                   # DDA solver (FFT-GMRES, convergence)
    mom2d.py                      # MoM 2D solver
    mom3d.py                      # MoM 3D solver (LU, GMRES, ACA)
    gen_dda.py                    # DDA batch runner (cache, convergence fit)
    gen_mom_2d.py                 # MoM 2D batch runner
    gen_mom_3d.py                 # MoM 3D batch runner
    gen_mesh_2d_rect.py           # 2D perimeter mesh generator
    gen_mesh_3d_rect.py           # 3D surface mesh generator

  examples/
    basic_usage.py                # Single-point demonstration
    run_mom3d_sweep.py            # 3D MoM sweep entry point
    run_mom2d_sweep.py            # 2D MoM sweep entry point
    run_dda_sweep.py              # DDA sweep entry point

  tests/
    test_vector.py                # V3, D3, Mesh unit tests
    test_dda_core.py              # DDA correctness & regression
    test_mom2d.py                 # 2D MoM mesh, solver, regression
    test_mom3d.py                 # 3D MoM mesh, solver, ACA, convergence

  data/                           # Caches & CSV output (auto-created)
    csvdata/
    _mom3d_cache/
    _mom2d_cache/
    _dda_cache/

  docs/                           # Documentation
    user_guide.pdf
\end{lstlisting}

\section{Prerequisites}

\begin{itemize}
    \item Python 3.9+ with NumPy and SciPy (installed automatically by
          \texttt{pip install -e .}).
    \item Cached results are stored in \texttt{data/\_mom3d\_cache/},
          \texttt{data/\_mom2d\_cache/}, and \texttt{data/\_dda\_cache/}.
          Delete these directories to force recomputation.
\end{itemize}

\section{Running the MoM 3D Sweep}

\subsection{Command-Line Options}

\begin{lstlisting}[language=bash]
python examples/run_mom3d_sweep.py [--er ER_LIST] [--solver direct|aca] [--dry-run]
\end{lstlisting}

\begin{center}
\begin{tabular}{@{}llp{6cm}@{}}
\toprule
Flag & Values & Description \\
\midrule
\texttt{--er} & Comma-separated floats
    & Relative permittivities to sweep.
      Default: \texttt{10,20,100,1e10} (where \texttt{1e10} $\approx$ PEC). \\
\texttt{--solver} & \texttt{direct} (default), \texttt{aca}
    & Solver backend.  \texttt{direct} uses LU factorization;
      \texttt{aca} uses ACA-compressed H-matrix + GMRES (recommended for
      large meshes, $N > 2000$). \\
\texttt{--dry-run} & (no value)
    & Print the shapes, $\varepsilon_r$ values, and $n_d$ sequences
      that would be computed, then exit. \\
\bottomrule
\end{tabular}
\end{center}

\subsection{Shape Grid}
\label{sec:mom3d_shapegrid}

The sweep covers 15 aspect ratios $(u,v) = (l_y/l_x,\; l_z/l_x)$:
%
\begin{quote}
\ttfamily
(1.0, 1.0), (1.0, 0.7), (1.0, 0.5), (1.0, 0.3), (1.0, 0.2),
(0.7, 0.7), (0.7, 0.5), (0.7, 0.3), (0.7, 0.2),
(0.5, 0.5), (0.5, 0.3), (0.5, 0.2),
(0.3, 0.3), (0.3, 0.2), (0.2, 0.2)
\end{quote}
%
with $l_x = 1$ fixed.

\subsection{Convergence}

For each shape, the solver runs 4--6 mesh resolutions $(n_d)$ chosen
adaptively by \texttt{adaptive\_nd()}.  A power-law fit
$P(n_d) = P_\infty + C\, n_d^{-\beta}$ extrapolates to the continuum
limit.  Individual solves and converged fits are cached on disk.

\subsection{Examples}

\begin{lstlisting}[language=bash]
# Preview what will be computed
python examples/run_mom3d_sweep.py --dry-run

# Sweep with default permittivities (direct LU)
python examples/run_mom3d_sweep.py

# Single permittivity with ACA solver
python examples/run_mom3d_sweep.py --er 10 --solver aca

# Custom permittivity range
python examples/run_mom3d_sweep.py --er 1.5,3,6,12,100
\end{lstlisting}

\subsection{Output}

Results are printed to stdout as the sweep progresses.  The depolarization
factors $N_x, N_y, N_z$ and their sum $N_{\Sigma} = N_x + N_y + N_z$ are
shown for each shape--$\varepsilon_r$ pair.

For bulk data export, use the standalone \texttt{gen\_mom\_3d.main()}
function (run \texttt{python -c "from espol\_cuboids.gen\_mom\_3d import main; main()" --dry-run}).

\section{Running the MoM 2D Sweep}

\subsection{Command-Line Options}

\begin{lstlisting}[language=bash]
python examples/run_mom2d_sweep.py [--er ER_LIST] [--dry-run]
\end{lstlisting}

\begin{center}
\begin{tabular}{@{}llp{6cm}@{}}
\toprule
Flag & Values & Description \\
\midrule
\texttt{--er} & Comma-separated floats
    & Relative permittivities.  Default: \texttt{0.7,1.5,3,6,1e3,1e5}. \\
\texttt{--dry-run} & (no value)
    & Print the shapes and $n_d$ sequences that would be computed, then exit. \\
\bottomrule
\end{tabular}
\end{center}

The 2D solver uses only the direct LU backend (no ACA/GMRES option), as
problem sizes remain small ($n \lesssim 500$).

\subsection{Shape Grid}

23 aspect ratios $r = l_y / l_x$ with $l_x = 1$:
%
\begin{quote}
\ttfamily
0.1, 0.15, 0.2, 0.3, 0.35, 0.4, 0.5, 0.6, 0.667, 0.8, 0.9,
1.0, 1.1, 1.25, 1.5, 1.75, 2.0, 2.5, 3.0, 3.33, 5.0, 7.0, 10.0
\end{quote}

\subsection{Convergence}

$n_d$ sequences are chosen adaptively.  A power-law fit extrapolates to
$n_d \to \infty$.  Individual solves and fits are cached in
\texttt{data/\_mom2d\_cache/}.

\subsection{Examples}

\begin{lstlisting}[language=bash]
# Preview
python examples/run_mom2d_sweep.py --dry-run

# Full sweep
python examples/run_mom2d_sweep.py

# Custom permittivities
python examples/run_mom2d_sweep.py --er 2,5,10
\end{lstlisting}

\subsection{Output}

For each shape--$\varepsilon_r$ pair, the normalized depolarization factors
\[
N_x = \frac{1}{P_{xx}/A}, \qquad
N_y = \frac{1}{P_{yy}/A}, \qquad
N_{\Sigma} = N_x + N_y
\]
are printed to stdout.  Results are cached on disk for reuse.

\section{Running the DDA Sweep}

\subsection{Command-Line Options}

\begin{lstlisting}[language=bash]
python examples/run_dda_sweep.py --mode fixed|converge [--er ER_LIST] [--dry-run]
\end{lstlisting}

\begin{center}
\begin{tabular}{@{}llp{5.5cm}@{}}
\toprule
Flag & Values & Description \\
\midrule
\texttt{--mode} & \texttt{fixed}, \texttt{converge}
    & \texttt{fixed}: single lattice spacing $d = s_{\min}/10$.
      \texttt{converge}: multiple $d$ levels with $d \to 0$ fit. \\
\texttt{--er} & Comma-separated floats
    & Relative permittivities.
      Default: \texttt{0.5,1.25,1.5,1.75,2,3,5,8,10,12.5}. \\
\texttt{--dry-run} & (no value)
    & Print the shapes and lattice parameters, then exit. \\
\bottomrule
\end{tabular}
\end{center}

The DDA sweep shares the same 15-shape grid as the MoM 3D sweep
(Section~\ref{sec:mom3d_shapegrid}).  The FFT-accelerated GMRES backend is
used for all solves ($N \gtrsim 1000$).

\subsection{Modes}

\begin{description}
    \item[\texttt{fixed}] Each shape is solved at a single lattice spacing
        $d = \min(l_y, l_z) / 10$.  Fast, suitable for quick surveys.
        Output: \texttt{data/csvdata/dda3d.csv}.

    \item[\texttt{converge}] Each shape is solved at 4--5 values of $d$
        (chosen so that $N \le 20\,000$ dipoles per solve).  A power-law
        $P(d) = P_\infty + C\, d^\beta$ fit yields the continuum limit.
        Output:
        \texttt{data/csvdata/dda\_convergence\_raw.csv} (individual solves) and
        \texttt{data/csvdata/dda\_convergence\_cnv.csv} (extrapolated values).
\end{description}

\subsection{Examples}

\begin{lstlisting}[language=bash]
# Preview fixed mode
python examples/run_dda_sweep.py --mode fixed --dry-run

# Quick fixed-d sweep (all er, all shapes)
python examples/run_dda_sweep.py --mode fixed

# Converge mode for er=10 only
python examples/run_dda_sweep.py --mode converge --er 10

# Converge mode with multiple er values
python examples/run_dda_sweep.py --mode converge --er 2,5,10
\end{lstlisting}

\subsection{Output}

Depolarization factors $N_x, N_y, N_z$ and their sum are printed per
shape--$\varepsilon_r$--$d$ combination.  CSV files are written to the
\texttt{data/csvdata/} directory.

\section{Interpreting Results}

\subsection{Depolarization Factors}

The depolarization factor $N_i$ (for axis $i$) relates the internal field to
the applied field in an ellipsoid:
\[
E_i^{\rm int} = E_i^{\rm app} - \frac{N_i}{\varepsilon_0} P_i
\]
For cuboids, $N_i$ is shape- and permittivity-dependent (unlike ellipsoids,
where it depends only on shape).  The sum rule
\[
N_x + N_y + N_z = 1 \qquad \text{(3D)},\qquad
N_x + N_y \approx 1 \qquad \text{(2D)}
\]
holds exactly for ellipsoids and approximately for cuboids at moderate
$\varepsilon_r$.

\subsection{Convergence Extrapolation}

The power-law fit $P(d) = P_\infty + C\,d^\beta$ is justified by the leading
discretization-error term for smooth geometries.  The exponent $\beta$ is
typically 1--2 for the MoM and 1--3 for the DDA.  Fits are monitored;
if the residual is large or $\beta$ falls outside expected bounds,
the finest-resolution value is used as a fallback.

\subsection{PEC Limit}

For the DDA, the PEC limit is obtained by fitting
$\alpha(\varepsilon_r) = \alpha_{\rm PEC} + A/\varepsilon_r$.
For the MoM, $\varepsilon_r = 10^{10}$ is used as a proxy (the
system matrix becomes nearly singular for true PEC, but
$\varepsilon_r = 10^{10}$ gives the correct limit to $\sim$4 significant
digits).

\section{Caching}

All three sweeps cache individual solves on disk:

\begin{itemize}
    \item \textbf{MoM 3D:} \texttt{data/\_mom3d\_cache/} — keyed by
          $\varepsilon_r$, shape, $n_d$, $\beta$, and solver type
          (direct/ACA).
    \item \textbf{MoM 2D:} \texttt{data/\_mom2d\_cache/} — keyed by
          $\varepsilon_r$, shape, $n_d$, $\beta$.
    \item \textbf{DDA:} \texttt{data/\_dda\_cache/} — keyed by
          $\varepsilon_r$, shape, and lattice spacing $d$.
\end{itemize}

Cache keys include version tags; changing the implementation version
automatically invalidates stale cache entries.  To force a full
recomputation, delete the relevant cache directory.

\section{Running Tests}

The test suite is run with pytest from the \texttt{espol-cuboids/} directory:

\begin{lstlisting}[language=bash]
python -m pytest tests/ -v                  # all 54 tests
python -m pytest tests/ -v -m "not slow"    # skip full MoM/DDA solves
\end{lstlisting}

Tests cover vector operations, mesh generation, MoM 2D/3D solvers, the
DDA solver, the ACA backend, convergence extrapolation, and regression
against known values.

\end{document}
